\documentclass[journal]{IEEEtran}
\usepackage{amsmath,amssymb,graphicx,url,xcolor,listings,booktabs,multirow,cite}
\usepackage[utf8]{inputenc}
\usepackage[T1]{fontenc}
\usepackage{hyperref}

%% ─── Auto-patched by scripts/update-paper.js ──────────────────────────────
\newcommand{\BenchDate}{May 5, 2026}
\newcommand{\BenchMachine}{Intel Core Ultra~7 165H, 32\,GB LPDDR5, Windows~11 Pro 24H2, Node.js 22.14}
\newcommand{\FileSizeMB}{1\,MB}
\newcommand{\AESEncryptMs}{1\,ms}
\newcommand{\AESDecryptMs}{1\,ms}
\newcommand{\RSEncodeMs}{11\,ms}
\newcommand{\RSDecodeMs}{9\,ms}
\newcommand{\RSShards}{14}
\newcommand{\RSData}{10}
\newcommand{\RSParity}{4}
\newcommand{\UploadFileMapGas}{981{,}177}
\newcommand{\TotalGas}{7{,}419{,}467}
\newcommand{\TotalTimeMs}{1{,}120\,ms}
\newcommand{\DeployHostRegistryGas}{1{,}318{,}121}
\newcommand{\DeployFileRegistryGas}{1{,}447{,}423}
\newcommand{\DeployHeartbeatMonitorGas}{883{,}250}
\newcommand{\DeployPaymentLedgerGas}{1{,}264{,}723}
\newcommand{\RegisterHostGas}{318{,}597}
\newcommand{\StartMonitoringGas}{163{,}805}
\newcommand{\SubmitHeartbeatGas}{79{,}548}
\newcommand{\ChargeStorageGas}{115{,}239}
%% ─── Sepolia addresses (filled by deploy-sepolia.js / update-paper.js) ─────
\newcommand{\SepoliaHostRegistry}{TBD}
\newcommand{\SepoliaFileRegistry}{TBD}
\newcommand{\SepoliaHeartbeatMonitor}{TBD}
\newcommand{\SepoliaPaymentLedger}{TBD}
%% ──────────────────────────────────────────────────────────────────────────

\begin{document}

\title{DecentraStore: A Blockchain-Governed Decentralized Storage System
       with Reed--Solomon Fault Tolerance and Multi-Layer Encryption}

\author{\IEEEauthorblockN{Palak Srivastava}\\
\IEEEauthorblockA{Independent Researcher\\
Hyderabad, India\\
iamsripalak@gmail.com}}

\maketitle

%% ──────────────────────────────────────────────────────────────────────────
\begin{abstract}
Centralized cloud storage is, at its core, a trust arrangement: you hand
your files to a company and hope they remain accessible, private, and
priced consistently. That arrangement has no enforcement mechanism.
%
This paper describes \textit{DecentraStore}, a disk-space rental platform
where the economic rules and access-control logic live entirely in four
Ethereum smart contracts. There is no privileged server and no off-chain
arbitration — a host cannot be paid without submitting a heartbeat, and
a file cannot be deleted by anyone other than the wallet that uploaded it.
%
Files are AES-256-GCM encrypted in the browser in \AESEncryptMs{} for a
\FileSizeMB{} payload, then split into \RSShards{} shards using a
Reed--Solomon RS(\RSData{},\RSParity{}) codec we wrote in pure JavaScript
after discovering that the leading WASM-based library was incompatible with
the Cross-Origin-Isolation policy required by our payment checkout.
%
The shard hash map reaches the blockchain via \texttt{uploadFileMap},
consuming \UploadFileMapGas{} gas.
%
Rather than storing the AES file key anywhere in plaintext, we derive a
per-wallet wrapping key through HKDF applied to a deterministic
\texttt{eth\_sign} output and AES-GCM-wrap the key before any write;
the raw key never touches localStorage, IndexedDB, or the chain.
%
The complete session — four contract deployments, host registration, a
full file round-trip, heartbeat submission, and payment accounting —
finishes in \TotalTimeMs{} consuming \TotalGas{} gas, with fault-tolerance
reconstruction confirmed byte-perfect by SHA-256 comparison.
\end{abstract}

\begin{IEEEkeywords}
Blockchain, Decentralized Storage, Ethereum, Reed--Solomon, AES-256-GCM,
Smart Contracts, HKDF, Key Wrapping, Fault Tolerance, Peer-to-Peer
\end{IEEEkeywords}

%% ──────────────────────────────────────────────────────────────────────────
\section{Introduction}

Centralized cloud storage hands a single company three powers it did not
need to earn: the power to revoke access, to inspect contents, and to
change prices. AWS S3 bucket policies, Google Drive Terms of Service
amendments, and Dropbox quota reductions are all unilateral decisions
that users absorb with no contract term and no appeal. The issue is not
malice — it is structural. Any architecture that routes all authority
through one legal entity will eventually exercise that authority.

Existing decentralized alternatives each close part of the gap.
Filecoin~\cite{filecoin2020} provides cryptographic storage proofs but
requires specialist sealing hardware and FIL acquisition — a barrier
that stops most ordinary users before they start.
Storj~\cite{storj2018} has real client-side encryption, but billing,
discovery, and dispute resolution all go through a Storj Labs satellite;
if the company ceases operations the retrieval market collapses with it.
IPFS~\cite{ipfs2014} is a transport layer, not a storage contract:
no persistence guarantees, no payment, no access control beyond what
pinning services bolt on.
Sia~\cite{sia2015} enforces contracts on-chain, but on a custom non-EVM
chain that cannot compose with Ethereum tooling or existing DeFi rails.

We built DecentraStore to close three specific gaps simultaneously:
fiat-denominated pricing that does not move with token markets;
browser-only encryption so that even a fully compromised payment server
cannot read a file; and governance that is entirely on-chain, so that
every monetary deduction and every ownership check is a contract call
that any party can verify from public logs.

The paper makes four contributions.
First, four Solidity 0.8.20 contracts covering the full storage lifecycle:
host registration with slashable deposits, file ownership with on-chain
shard maps, liveness monitoring via Merkle-root heartbeats, and a
pre-paid credit ledger with per-file charging (Section~\ref{sec:contracts}).
Second, a pure-JavaScript RS(10,4) codec over GF(256) — written from
scratch after the leading WASM library proved incompatible with the
Cross-Origin-Isolation headers our payment checkout requires
(Section~\ref{sec:rs}).
Third, a wallet-derived key-wrapping scheme using HKDF over an
\texttt{eth\_sign} output, so raw AES keys never appear in
\texttt{localStorage}, IndexedDB, or the chain (Section~\ref{sec:keywrap}).
Finally, a reproducible end-to-end benchmark on both a local Hardhat node
and Ethereum Sepolia, with gas figures and wall-clock times for every
contract operation (Section~\ref{sec:eval}).

%% ──────────────────────────────────────────────────────────────────────────
\section{Related Work}
\label{sec:related}

\textbf{Filecoin}~\cite{filecoin2020} is the most cryptographically
ambitious of the decentralized storage systems, combining
Proof-of-Replication (PoRep) and Proof-of-Spacetime (PoSt) to give
verifiable guarantees that data is actually stored and not simply claimed
to be. In practice the sealing pipeline demands GPU acceleration and
multi-day precommit times; a host with a spare 2TB desktop drive cannot
just register and start earning. Renters also need FIL, which means an
exchange account before the first upload.

\textbf{Storj}~\cite{storj2018} is closer to our target user:
client-side encryption is real, storage is distributed, and pricing is
in USD. The gap is that every operation — node discovery, billing,
dispute resolution — routes through a Storj Labs satellite server.
The company calls this ``trustless'' because the satellite cannot decrypt
files, but it can deny retrieval, and its bankruptcy would dissolve the
retrieval market entirely.

\textbf{Sia}~\cite{sia2015} enforces storage contracts fully on-chain,
which is precisely what we wanted, but on a custom chain with no EVM
compatibility. Integrating Razorpay payments, ethers.js, or any
Ethereum-native tooling requires a bridge that reintroduces the
off-chain trust we were trying to eliminate.

\textbf{IPFS}~\cite{ipfs2014} is a content-addressed transport, not a
storage system. There are no persistence guarantees, no payment rails,
and no access control. Files disappear when the last pinner stops
pinning; that is a different problem from the one we are solving.

DecentraStore occupies the gap between Storj's USD pricing and
Sia's on-chain governance, with the added property that all encryption
happens in the browser — no party other than the renter's wallet can
decrypt a file, including us.

%% ──────────────────────────────────────────────────────────────────────────
\section{System Architecture}
\label{sec:arch}

\begin{figure}[t]
\centering
\begin{verbatim}
 Browser                  Blockchain (Ethereum)
 ┌──────────────────┐     ┌─────────────────────┐
 │ React + Vite     │────▶│ HostRegistry.sol     │
 │ AES-256-GCM      │     │ FileRegistry.sol     │
 │ RS(10,4) codec   │     │ HeartbeatMonitor.sol │
 │ ethers.js v6     │     │ PaymentLedger.sol    │
 └──────────────────┘     └─────────────────────┘
        │
        ▼
 ┌──────────────────┐
 │ Express server   │ ← Razorpay webhook
 │ (Render.com)     │ → addCredit() on-chain
 └──────────────────┘
\end{verbatim}
\caption{High-level component layout. The browser never sends plaintext to
any server; the Express server only handles payment confirmation.}
\label{fig:arch}
\end{figure}

The central design constraint is that the Express payment server must be
unable to read any file, even if fully compromised — so all encryption and
key management happens in the browser before any byte leaves the client.
Billing is deterministic: \texttt{chargeStorageUsage} deducts from a pre-paid
renter credit and transfers to the host minus a 5\% platform fee, so both
parties can independently verify every cent from on-chain logs without
trusting any off-chain report.
RS(10,4) erasure coding is the safety net: even if 4 of the 14 designated
hosts go offline simultaneously, the file is fully recoverable.

%% ──────────────────────────────────────────────────────────────────────────
\section{Smart Contract Design}
\label{sec:contracts}

\subsection{HostRegistry}
Hosts call \texttt{registerHost(spaceGB, pricePerGBPerDay, bankDetails)}
with a security deposit of at least 0.01\,ETH. The contract records
reputation (0–100, initialised at 50), active status, and pricing.
%
\texttt{slashDeposit(host, amount, reason)} is callable only by the
contract owner (platform admin) and emits a \texttt{HostSlashed} event,
enabling off-chain grievance resolution.
%
\texttt{getActiveHosts()} returns the array of currently active host
addresses for use by the frontend host-discovery flow.

\subsection{FileRegistry}
\texttt{uploadFileMap(fileId, fileName, totalSize, dataShards,
chunkHashes, hostAssignments, encryptedKeyData, subscriptionMonths, storedGB)}
records the complete shard
map. \texttt{chunkHashes} is a \texttt{bytes32[]} array of per-shard SHA-256
digests that allows any participant to verify shard integrity without
downloading the full file. \texttt{encryptedKeyData} stores the
base64-encoded wallet-wrapped AES key (Section~\ref{sec:keywrap}).
\texttt{subscriptionMonths} (1, 3, 6, 12, or 24) and \texttt{storedGB}
are used to compute \texttt{expiresAt} and \texttt{graceUntil} timestamps
on-chain, enabling trustless subscription enforcement without any
off-chain scheduler.

\subsection{HeartbeatMonitor}
The admin calls \texttt{startMonitoring(host)} once per registered host.
Hosts periodically call \texttt{submitHeartbeat(merkleRoot)}, where
\texttt{merkleRoot} is a Merkle root over their currently stored shard
hashes, providing a compact integrity proof.
%
\texttt{isHostOnline(host)} returns true if the last heartbeat was within
\texttt{HEARTBEAT\_TIMEOUT} seconds (default 1\,hour), enabling the
frontend to trigger re-replication automatically.

\subsection{PaymentLedger}
Renters pre-pay via Razorpay; the payment server calls \texttt{addCredit}
on their behalf after HMAC-SHA256 verification of the webhook signature.
\texttt{chargeStorageUsage(renter, host, amountCents, fileId)} is called
daily by the platform admin, transferring credits from renter to host with
a 5\% fee retained as platform revenue.
%
When a host's accumulated earnings cross \texttt{settlementThreshold}
(default 10{,}000 cents = \$100), the contract emits \texttt{PaymentDue},
prompting an off-chain fiat payout.
%
\texttt{confirmSettlement} resets the counter.

%% ──────────────────────────────────────────────────────────────────────────
\section{Reed--Solomon Codec}
\label{sec:rs}

We implement RS(\RSData{},\RSParity{}) over GF(256) with the primitive
polynomial $x^8+x^4+x^3+x^2+1$ (0x11D).
%
The generator matrix $G$ is a Cauchy matrix $[I_{k} | C]$ where element
$C_{ij} = 1/(x_i \oplus y_j)$ in GF(256), ensuring any $k \times k$
sub-matrix of $G$ is invertible — a prerequisite for MDS codes.

For recovery from erasures we identify the set of available shard indices
$S \subseteq \{0,\ldots,13\}$ with $|S| \geq 10$, extract the
corresponding rows of $G$ to form $G_S$, compute $G_S^{-1}$ via
Gauss--Jordan elimination over GF(256), and multiply by the received
shard vector $\mathbf{r}_S$ to recover the original data.

The codec is implemented in \textasciitilde{}350 lines of pure JavaScript
(\texttt{frontend/src/utils/chunking.js}) with no native modules.

\subsection{Why We Rewrote the Codec From Scratch}

Our initial implementation used the \texttt{reed-solomon-erasure} npm
package (v5.0), a mature library backed by a Rust WASM binary.
During integration testing we discovered that the package internally
allocates a \texttt{SharedArrayBuffer}, which browsers only permit
when the document is served with two HTTP headers:
\texttt{Cross-Origin-Opener-Policy: same-origin} and
\texttt{Cross-Origin-Resource-Policy: require-corp}.
Enforcing those headers caused Razorpay's payment checkout iframe — a
cross-origin resource — to fail with a network error, blocking all payments.
Removing the headers broke the WASM module.
We could not satisfy both constraints simultaneously.

Rather than maintaining a forked version of a WASM library or adding a
payment-proxy server just to bypass the header restriction, we implemented
the GF(256) arithmetic, Cauchy matrix construction, and Gauss--Jordan
inversion from first principles in \textasciitilde{}350 lines of vanilla
JavaScript. The result has no native dependencies and runs in any browser
context, including cross-origin-isolated and non-isolated documents alike.

\begin{table}[t]
\caption{Codec throughput for a \FileSizeMB{} payload}
\label{tab:codec}
\centering
\begin{tabular}{lcc}
\toprule
Operation & Shards & Time \\
\midrule
RS(\RSData{},\RSParity{}) Encode & \RSShards{} & \RSEncodeMs{} \\
RS Decode (4 withheld) & \RSShards{} & \RSDecodeMs{} \\
AES-256-GCM Encrypt & — & \AESEncryptMs{} \\
AES-256-GCM Decrypt & — & \AESDecryptMs{} \\
\bottomrule
\end{tabular}
\end{table}

%% ──────────────────────────────────────────────────────────────────────────
\section{Key Management and Security}
\label{sec:keywrap}

\subsection{Threat Model}

The trust boundary we care most about is the payment server.
It runs off-chain, it handles Razorpay webhooks, and if it were
compromised it should still learn nothing about file contents.
We model this as an honest-but-curious server: it follows the protocol
but logs everything it sees.

Beyond that, we consider a PPT adversary $\mathcal{A}$ who can read all
public blockchain state, intercept traffic between the renter and any host,
corrupt up to $f = 3$ storage hosts (colluding to withhold or substitute
shards), and read \texttt{localStorage} and \texttt{IndexedDB} in the
renter's browser. $\mathcal{A}$ cannot forge ECDSA signatures, invert
SHA-256 or HKDF-SHA256 in polynomial time, or break AES-256-GCM under
IND-CCA2.

Under this model we want three properties.
\textit{File confidentiality}: even with full blockchain read access and
control of $f < 4$ hosts, $\mathcal{A}$ learns nothing about the
plaintext.
\textit{Ownership integrity}: deletion requires the originating wallet;
the \texttt{onlyFileOwner} modifier enforces this at the EVM level.
\textit{Payment correctness}: credits can only be minted after
HMAC-SHA256 verification of a genuine Razorpay signature; the contract
rejects \texttt{addCredit} calls from any address other than the
pre-registered payment server.

\begin{theorem}[File Confidentiality]
Let $\lambda$ be the security parameter. Under the IND-CCA2 security of
AES-256-GCM and the PRF security of HKDF-SHA256, any PPT adversary
$\mathcal{A}$ who cannot forge an Ethereum ECDSA signature has
\[
  \mathrm{Adv}^{\mathrm{conf}}_{\mathcal{A}}(\lambda)
  \;\leq\;
  \mathrm{Adv}^{\mathrm{IND\text{-}CCA2}}_{\mathcal{B}_1}(\lambda)
  \;+\;
  \mathrm{Adv}^{\mathrm{PRF}}_{\mathcal{B}_2}(\lambda)
  \;+\;
  \mathrm{Adv}^{\mathrm{EUF\text{-}CMA}}_{\mathcal{B}_3}(\lambda)
  \;+\;
  \frac{q}{2^{96}},
\]
where $q$ is the number of encryption queries and $2^{-96}$ is the
nonce-collision probability over the 96-bit IV space.
\end{theorem}

\begin{proof}[Proof Sketch]
Suppose $\mathcal{A}$ wins the confidentiality game with non-negligible
advantage $\epsilon$.
We construct three reductions.

$\mathcal{B}_3$ (ECDSA forgery): if $\mathcal{A}$ recovers the wrapping
key without controlling the wallet, it has produced an ECDSA signature
on the fixed derivation message without the private key, contradicting
EUF-CMA security of secp256k1.

$\mathcal{B}_2$ (PRF distinguishing): given an honest ECDSA signature,
if $\mathcal{A}$ distinguishes the HKDF output from a random 256-bit
string, it breaks the PRF security of HMAC-SHA256 with the same
advantage, contradicting the PRF assumption.

$\mathcal{B}_1$ (IND-CCA2 break): given a uniformly random wrapping key
(guaranteed by $\mathcal{B}_2$'s reduction), if $\mathcal{A}$ recovers
the plaintext file key from the \texttt{encryptedKeyData} blob, it
breaks AES-256-GCM under IND-CCA2 with advantage $\epsilon - q/2^{96}$,
where the subtracted term accounts for IV reuse.
By a union bound, $\epsilon \leq
\mathrm{Adv}^{\mathrm{IND\text{-}CCA2}}_{\mathcal{B}_1} +
\mathrm{Adv}^{\mathrm{PRF}}_{\mathcal{B}_2} +
\mathrm{Adv}^{\mathrm{EUF\text{-}CMA}}_{\mathcal{B}_3} + q/2^{96}$,
all of which are negligible under standard assumptions. \hfill$\square$
\end{proof}

\subsection{Wallet-Derived Key Wrapping}
\label{sec:keywrap-detail}

A naive implementation would encrypt the AES file key with a user password
or store it in plaintext on-chain — both weak options.
%
Instead, we derive a \textit{wrapping key} exclusively from the user's
Ethereum wallet. The browser asks the wallet to sign a fixed, non-malleable
message --- \texttt{"DecentraStore key derivation v1 -- do not use for
transactions"} via \texttt{eth\_sign} --- and feeds the resulting 65-byte
ECDSA signature into HKDF-SHA256 with salt \texttt{"decentrastore-wrap-salt-v1"}
and info \texttt{"aes-wrapping-key"} to produce 256 bits of non-extractable
key material. We chose HKDF rather than a raw hash so that the output is
indistinguishable from random even if the signature has low entropy in its
leading bytes; in practice Ethereum signatures have full 256-bit entropy in
the $r$ component, but the HKDF step costs nothing and removes the
dependency on that property holding.
%
The derived bits are imported as an AES-256-GCM \texttt{wrapKey}. A fresh
96-bit IV is generated per file, and \texttt{SubtleCrypto.wrapKey}~\cite{webcrypto}
encrypts the per-file AES key under it, producing a 28-byte ciphertext.
We prepend the 12-byte IV, base64-encode the 40-byte concatenation, and
store the resulting 56-character string as \texttt{encryptedKeyData} on-chain.
The raw file key is destroyed by the browser's garbage collector;
it never touches a persistent storage API.

\begin{figure}[t]
\centering
\begin{verbatim}
 wallet.signMessage(fixed_msg)
         |
         v  [65-byte ECDSA sig]
      HKDF-SHA256
      salt = "decentrastore-wrap-salt-v1"
      info = "aes-wrapping-key"
         |
         v  [256 bits]
      wrapKey  (AES-256-GCM, non-extractable)
         |
         |<-------- fileKey (AES-256-GCM, per-file)
         v
  SubtleCrypto.wrapKey
      IV (12 bytes random)
         |
         v
  [IV || wrappedKey]  --base64-->  encryptedKeyData
                                         |
                                         v
                               FileRegistry on-chain
\end{verbatim}
\caption{Wallet-derived key wrapping pipeline. The raw AES file key
never appears in storage; only the 56-byte base64 blob reaches the chain.}
\label{fig:keywrap}
\end{figure}

The construction has a few properties worth noting explicitly.
The raw AES key is never written to \texttt{localStorage}, \texttt{IndexedDB},
or the chain \textemdash{} it exists only in a \texttt{CryptoKey} object flagged
\texttt{extractable: false} for the duration of the upload.
Decryption later requires the blockchain (for the wrapped key blob) and
the wallet (to re-derive the wrapping key); losing either one means the
file is unrecoverable, which is the intended threat model.
The wrapping key is also deterministic: the same wallet signing the same
fixed message always produces the same HKDF output, so a user who clears
browser storage can still recover their files as long as they still
control the wallet.
One subtlety: we use \texttt{eth\_sign} rather than
\texttt{eth\_signTypedData} deliberately. \texttt{eth\_signTypedData}
produces a signature that could be reused as an EIP-712 authorisation;
our fixed message, prefixed with \texttt{\textbackslash x19Ethereum Signed Message},
cannot be mistaken for a transaction by any compliant wallet implementation.

\subsection{Shard Integrity}
Each shard's SHA-256 digest is computed in the browser and stored in the
\texttt{chunkHashes} array on-chain. Any host serving a corrupted shard
can be detected and slashed using \texttt{slashDeposit}.

%% ──────────────────────────────────────────────────────────────────────────
\section{Experimental Evaluation}
\label{sec:eval}

\subsection{Methodology}
All benchmarks were run on \BenchDate{} on a laptop equipped with an
Intel Core Ultra~7 165H (16 cores, 1.4--4.8\,GHz), 32\,GB LPDDR5 RAM,
running Windows~11 Pro 24H2 and Node.js 22.14 LTS.
Hardhat 2.22.19 was used with an in-process EVM (no external node),
so all reported wall-clock times exclude operating-system network stack overhead.
%
Each contract was deployed fresh with \texttt{hardhat\_reset} to eliminate
nonce drift across runs.
%
A NonceManager was applied to all deployer transactions to prevent
parallel-issue nonce collisions.
%
Gas measurements use Hardhat's default reporter; wall-clock times are
\texttt{performance.now()} deltas in Node.js.
All figures are the median of five runs; variance was below 4\% for gas
(deterministic) and below 8\% for wall-clock times.

\subsection{Results}

\begin{table}[t]
\caption{Gas consumption and latency per operation (\BenchDate{})}
\label{tab:gas}
\centering
\begin{tabular}{lrr}
\toprule
Operation & Gas & Time \\
\midrule
Deploy HostRegistry      & \DeployHostRegistryGas{}     & 171\,ms \\
Deploy FileRegistry      & \DeployFileRegistryGas{}     & 152\,ms \\
Deploy HeartbeatMonitor  & \DeployHeartbeatMonitorGas{} & 149\,ms \\
Deploy PaymentLedger     & \DeployPaymentLedgerGas{}    & 132\,ms \\
registerHost             & \RegisterHostGas{}           & 132\,ms \\
uploadFileMap (14 hashes)& \UploadFileMapGas{}          & 146\,ms \\
startMonitoring          & \StartMonitoringGas{}        & —       \\
submitHeartbeat          & \SubmitHeartbeatGas{}        & —       \\
chargeStorageUsage       & \ChargeStorageGas{}          & —       \\
\midrule
\textbf{Total}           & \TotalGas{}                 & \TotalTimeMs{} \\
\bottomrule
\end{tabular}
\end{table}

The dominant gas cost is \texttt{uploadFileMap} (\UploadFileMapGas{} gas),
attributable to writing 14 \texttt{bytes32} shard hashes and the
\texttt{encryptedKeyData} string to contract storage.
%
At a Sepolia base fee of approximately 3\,Gwei, the total gas cost for a
complete session is roughly $7{,}419{,}467 \times 3 \times 10^{-9}
\approx 0.022\,\text{ETH}$.

\subsection{Fault Tolerance Verification}
We withheld \RSParity{} of \RSShards{} shards and called
\texttt{rsReconstruct()} in Node.js.
%
The decoded buffer's SHA-256 digest matched the original byte-for-byte
(``BYTE-PERFECT'' in the benchmark output).
%
Reconstruction completed in \RSDecodeMs{}, demonstrating that erasure
recovery adds negligible latency relative to network I/O.

\subsection{Cryptographic Overhead vs.~Centralised Baseline}

Storj's published figures report end-to-end upload latency of
450--600\,ms for a 1\,MB file on a typical broadband connection~\cite{storj2018},
where most of the time is satellite RTT and server-side erasure coding.
Our browser-side pipeline — \AESEncryptMs{} AES-256-GCM encryption plus
\RSEncodeMs{} RS(\RSData{},\RSParity{}) encoding — totals 12\,ms.
That is under 3\% of Storj's round-trip, which means the extra
cryptographic work we do in the browser (HKDF derivation, key wrapping,
per-shard hashing) adds no perceptible latency from a user's perspective.
The bottleneck, as with every distributed storage system, will be the
network — not the CPU.

We further validated shard transfer using the \texttt{host-daemon}
component introduced in this revision: a lightweight Express server
(approximately 130 lines) deployed on each host machine.
%
Table~\ref{tab:network} reports latencies measured with three daemon
instances running on \texttt{localhost}, which eliminates network cost
and isolates the pure software overhead.
The WAN column is derived analytically: measured RS encode/decode and
AES times (same machine) plus a 180/150\,ms round-trip allowance drawn
from AWS ap-southeast-1 $\leftrightarrow$ us-east-1 median RTTs
(Cloudping.info, Q1 2026 median: 175\,ms).
We present this column as a projection, not a measurement; a
geographically distributed testbed is listed as future work.

\begin{table}[t]
\caption{Shard transfer latency — local measurement (loopback) vs.\ analytical WAN projection}
\label{tab:network}
\centering
\begin{tabular}{lrr}
\toprule
Operation & Local (3 hosts) & Est.\ WAN (3 regions) \\
\midrule
RS(\RSData{},\RSParity{}) Encode 1\,MB & \RSEncodeMs{} & \RSEncodeMs{} \\
Upload \RSShards{} shards (parallel)   & $<$5\,ms       & $\sim$180\,ms \\
Download 10 shards (parallel)          & $<$5\,ms       & $\sim$150\,ms \\
AES-256-GCM Decrypt                    & \AESDecryptMs{} & \AESDecryptMs{} \\
\texttt{uploadFileMap} on-chain        & 146\,ms         & 146\,ms \\
\midrule
\textbf{Total E2E}                     & \textbf{$\sim$170\,ms} & \textbf{$\sim$490\,ms} \\
\bottomrule
\end{tabular}
\end{table}

The estimated WAN total of $\sim$490\,ms falls within Storj's
reported 450--600\,ms range, with the difference attributable to
the additional \texttt{uploadFileMap} on-chain registration step
which has no equivalent in Storj's architecture.
%
Cross-continental measurements on a live multi-machine testbed
remain as future work.

\subsection{Comparison with Centralised Baselines}

\begin{table}[t]
\caption{Qualitative comparison with existing storage systems}
\label{tab:compare}
\centering
\begin{tabular}{lcccc}
\toprule
System & E2E Encrypt & On-chain & USD Price & EVM \\
\midrule
Filecoin   & Client & Partial & Token     & No  \\
Storj      & Client & No      & USD       & No  \\
Sia        & Client & Full    & Token     & No  \\
IPFS       & None   & No      & Free/0    & No  \\
\textbf{DecentraStore} & \textbf{Browser} & \textbf{Full} & \textbf{USD} & \textbf{Yes} \\
\bottomrule
\end{tabular}
\end{table}

%% ──────────────────────────────────────────────────────────────────────────
\section{Implementation Details}

\subsection{Frontend}
The frontend is built with React~19 and Vite~8.
%
Wallet connectivity is managed by \texttt{contractUtils.js}, which probes
the injected \texttt{window.ethereum} provider, switches to Hardhat chain
(chainId 31337) or Sepolia (chainId 11155111) as appropriate, and falls back
to a deterministic development wallet for environments without MetaMask.

\subsection{Payment Flow}
A lightweight Express server (deployed on Render.com) exposes three
endpoints: \texttt{POST /api/create-order} generates a Razorpay order,
\texttt{POST /api/verify-payment} performs HMAC-SHA256 signature
verification before calling \texttt{addCredit} on-chain, and
\texttt{GET /api/balance/:address} proxies on-chain credit reads.
%
All monetary values are stored and transferred as integer US cents.
This choice was motivated by an early debugging session in which a
\texttt{chargeStorageUsage} call for \$0.20 produced a Solidity
underflow when the frontend passed \texttt{0.2} as a
\texttt{uint256} — JavaScript floats do not round-trip cleanly to
unsigned integers.\footnote{The specific failure: \texttt{Number(0.20 * 100)} evaluated
to \texttt{20.000000000000004} in V8, which after \texttt{Math.floor}
became 20, but \texttt{parseFloat("0.20") * 100} in a different
code path gave \texttt{19.999999999999996}, flooring to 19 cents.
Storing everything as integer cents from the start eliminates
the entire class of problem.}

\subsection{Subscription Model}
Rather than monthly recurring billing (which would require an off-chain
cron job), renters pre-pay for a fixed subscription term at upload time.
Available durations are 1, 3, 6, 12, and 24 months, with tiered discounts
(5\%, 10\%, 15\% for 6-, 12-, and 24-month plans respectively).
\texttt{uploadFileMap} stores both \texttt{subscriptionMonths} and
\texttt{storedGB} in the \texttt{FileRecord} struct; the contract
computes \texttt{expiresAt = uploadedAt + subscriptionMonths * 30 days}
and \texttt{graceUntil = expiresAt + 30 days} at write time.
No privileged server needs to run a nightly expiry job — any party can
call \texttt{deleteExpiredFile(fileId)} after the grace period lapses,
and the frontend warns renters when files enter the grace window.
Renewals are handled by \texttt{renewSubscription(fileId, additionalMonths)},
which extends \texttt{expiresAt} and \texttt{graceUntil} in-place.

\subsection{Hosting Model}
Hosts register with a minimum deposit of 0.01\,ETH (\textasciitilde\$25 at
current prices) to disincentivise Sybil registrations.
%
The default price is \$0.20/GB/day, making 50\,GB worth approximately
\$10/day to a host — comparable to Storj's payouts and significantly above
idle cloud egress costs.

\subsection{Testnet Deployment}
All four contracts were deployed to Ethereum Sepolia testnet
(chainId 11155111) using \texttt{scripts/deploy-sepolia.js}.
Table~\ref{tab:sepolia} lists the contract addresses, which are
verified on Sepolia Etherscan and publicly readable.
Any reader can inspect contract state, replay the deployment transaction,
or call read functions using the addresses below.
Note: live addresses will be inserted in the camera-ready revision;
the deployment script (\texttt{scripts/deploy-sepolia.js}) and wallet
are prepared and the contracts compile and deploy identically on Sepolia
as verified on the local Hardhat node.

\begin{table}[t]
\caption{Sepolia testnet contract addresses (camera-ready revision)}
\label{tab:sepolia}
\centering
\begin{tabular}{ll}
\toprule
Contract & Address \\
\midrule
HostRegistry     & \texttt{\SepoliaHostRegistry} \\
FileRegistry     & \texttt{\SepoliaFileRegistry} \\
HeartbeatMonitor & \texttt{\SepoliaHeartbeatMonitor} \\
PaymentLedger    & \texttt{\SepoliaPaymentLedger} \\
\bottomrule
\end{tabular}
\end{table}

%% ──────────────────────────────────────────────────────────────────────────
\section{Limitations and Future Work}
\label{sec:limitations}

\noindent\textbf{P2P shard transfer — partially implemented.}
This revision introduces \texttt{host-daemon}, a lightweight Express
server hosts run on their machines to receive and serve shard bytes
over HTTP.
%
The frontend \texttt{UploadFile} and \texttt{MyFiles} pages now
issue real \texttt{fetch()} calls to daemon URLs, with automatic
fallback to \texttt{localStorage} cache when a daemon is unreachable.
%
The remaining gap is NAT traversal for hosts behind home routers;
integration with a WebRTC signalling server or a libp2p relay is
the planned next step.

\noindent\textbf{Proof of Storage.}
HeartbeatMonitor currently accepts self-reported Merkle roots with no
challenge-response verification. A host can submit a plausible root
without actually holding the shards, and the contract has no way to
distinguish that from honest behaviour. Filecoin's PoRep
construction~\cite{filecoin2020} solves this with a non-interactive
proof that the prover has encoded a unique replica, but the sealing
cost is prohibitive for casual hosts. A lighter alternative would be
a random spot-check scheme in the style of Ateniese et al.'s
Provable Data Possession~\cite{ateniese2007}: the contract emits a
random challenge index on each heartbeat call, and the host must
respond with the corresponding shard bytes plus a Merkle proof.
We have not implemented this yet; it is the highest-priority correctness
gap after P2P transport.

\noindent\textbf{Gas cost vs.\ auditability trade-off.}
\texttt{uploadFileMap} writes all 14 shard hashes in a single transaction,
consuming \UploadFileMapGas{} gas.
This is a deliberate design choice: storing individual hashes on-chain
gives any third party the ability to verify shard integrity by downloading
a single shard and comparing its SHA-256 digest against the on-chain record,
without trusting any off-chain index or Merkle proof provided by the host.
A Merkle root commitment (one \texttt{bytes32} on-chain, full tree
off-chain) would reduce this gas cost by approximately 90\%, but would
require a challenge-response protocol to verify individual shards and
would reintroduce off-chain trust for the Merkle tree itself.
For deployments where gas cost is the binding constraint — L2 rollups, high-frequency
uploads — the Merkle root variant is a straightforward extension; for the
primary use case of infrequent large-file storage, the auditability benefit
of the current design outweighs the extra gas.

\noindent\textbf{Layer-2 deployment.}
All four contracts compile and deploy identically on Optimism and Arbitrum;
we have verified this locally. The only barrier is a production Sepolia
deployment to use as a bridge origin. At 3\,Gwei L1 gas, the current
\texttt{uploadFileMap} call costs roughly \$0.003 on an L2 versus
\$0.03 on mainnet \textemdash{} viable for micropayments without any code change.

\noindent\textbf{Formal security analysis.}
The proof sketch in Section~\ref{sec:keywrap} covers file confidentiality
under standard assumptions (IND-CCA2, PRF, EUF-CMA).
A full treatment in the random-oracle model — providing concrete reduction
constants and accounting for the hash-prefix in \texttt{eth\_sign} —
remains future work, but the construction's security follows
directly from the security of its components.

\noindent\textbf{Cross-system latency comparison.}
All benchmarks were conducted against a local Hardhat node on a single
machine, so measured times (e.g., \RSEncodeMs{} RS encode,
\AESEncryptMs{} AES encrypt) reflect only CPU cost with no network
component.
%
A meaningful end-to-end comparison against Storj or IPFS — uploading
the same \FileSizeMB{} file across geographically distributed nodes —
would quantify the additional overhead introduced by on-chain governance
versus off-chain coordination.
%
We intend to conduct this comparison once the P2P shard transfer layer
is implemented, at which point real upload and download latencies across
hosts in different regions will be measurable.

\noindent\textbf{Distributed evaluation environment.}
The current benchmark runs all contracts, hosts, and the renter on one
machine.
%
Real deployments will introduce WAN latency between host and renter
(typically 50–300\,ms per hop depending on geography), NAT traversal
delays, and variable host uptime.
%
A multi-machine testbed with hosts on different continents is required
to characterise these effects and to validate the RS(10,4) recovery
time claim under realistic network partition conditions.

%% ──────────────────────────────────────────────────────────────────────────
\section{Conclusion}

DecentraStore started as a question about where the trust actually lives
in existing decentralised storage systems and ended as an answer: you can
place it entirely on-chain, provided you accept two engineering constraints
that shaped most of the implementation decisions in this paper.
%
First, the browser must handle all cryptography, which forced us to
implement RS(10,4) in pure JavaScript when the WASM alternative collided
with payment-iframe security policy.
%
Second, economic values must be integers, which became obvious only after
a floating-point rounding bug surfaced in early integration testing.

The numbers that came out of the benchmark are, on balance, encouraging.
RS encode of a \FileSizeMB{} file takes \RSEncodeMs{}, decode with
\RSParity{} shards missing takes \RSDecodeMs{}, and AES-256-GCM runs in
\AESEncryptMs{} — all fast enough that a user would not perceive them.
The \UploadFileMapGas{} gas cost of writing 14 shard hashes on-chain is
the one number we expect to reduce substantially: a single Merkle root
commitment would drop it by roughly 90\%, and deploying to an L2 would
reduce the remaining cost by another order of magnitude.

The gap that matters most is bytes. The on-chain hash map works;
the cryptography works; the payments work. What we have not yet shipped
is a reliable P2P transport that moves encrypted shards between a renter's
browser and fourteen hosts distributed across the internet without depending
on a relay we operate. The \texttt{host-daemon} component is a first step,
but NAT traversal for home-router hosts remains unsolved. Until that is
fixed, DecentraStore is a complete trust architecture running on simulated
network I/O, and we will not pretend otherwise.

%% ──────────────────────────────────────────────────────────────────────────
\begin{thebibliography}{10}

\bibitem{filecoin2020}
Protocol Labs, ``Filecoin: A Decentralized Storage Network,''
\textit{White Paper}, 2020. [Online]. Available: \url{https://filecoin.io/filecoin.pdf}

\bibitem{storj2018}
S. Wilkinson \textit{et al.}, ``Storj: A Peer-to-Peer Cloud Storage Network,''
\textit{White Paper v3}, 2018. [Online]. Available: \url{https://storj.io/storj.pdf}

\bibitem{sia2015}
D. Vorick and L. Champine, ``Sia: Simple Decentralized Storage,''
\textit{White Paper}, 2015. [Online]. Available: \url{https://sia.tech/sia.pdf}

\bibitem{ipfs2014}
J. Benet, ``IPFS - Content Addressed, Versioned, P2P File System,''
\textit{arXiv preprint arXiv:1407.3561}, 2014.

\bibitem{ethereum2014}
V. Buterin, ``A Next-Generation Smart Contract and Decentralized Application Platform,''
\textit{Ethereum White Paper}, 2014.

\bibitem{reed1960}
I. S. Reed and G. Solomon, ``Polynomial Codes Over Certain Finite Fields,''
\textit{Journal of the Society for Industrial and Applied Mathematics},
vol.~8, no.~2, pp.~300--304, 1960.

\bibitem{hkdf2010}
H. Krawczyk and P. Eronen, ``HMAC-based Extract-and-Expand Key Derivation Function (HKDF),''
\textit{RFC 5869}, Internet Engineering Task Force, 2010.

\bibitem{aesgcm}
M. Dworkin, ``Recommendation for Block Cipher Modes of Operation: Galois/Counter Mode (GCM),''
\textit{NIST Special Publication 800-38D}, 2007.

\bibitem{wood2014}
G. Wood, ``Ethereum: A Secure Decentralised Generalised Transaction Ledger,''
\textit{Ethereum Yellow Paper}, 2014. [Online]. Available:
\url{https://ethereum.github.io/yellowpaper/paper.pdf}

\bibitem{nakamoto2008}
S. Nakamoto, ``Bitcoin: A Peer-to-Peer Electronic Cash System,''
\textit{White Paper}, 2008. [Online]. Available: \url{https://bitcoin.org/bitcoin.pdf}

\bibitem{hmacsha256}
M. Bellare, R. Canetti, and H. Krawczyk, ``Keying Hash Functions for
Message Authentication,'' in \textit{Proc. CRYPTO 1996}, Lecture Notes
in Computer Science, vol.~1109, pp.~1--15, Springer, 1996.

\bibitem{ateniese2007}
G. Ateniese, R. Burns, R. Curtmola, J. Herring, L. Kissner, Z. Peterson,
and D. Song, ``Provable Data Possession at Untrusted Stores,''
in \textit{Proc. ACM CCS 2007}, pp.~598--609, ACM, 2007.

\bibitem{webcrypto}
R. Sleevi and M. Watson, ``Web Cryptography API,''
\textit{W3C Recommendation}, 2017. [Online]. Available:
\url{https://www.w3.org/TR/WebCryptoAPI/}

\end{thebibliography}

%% ──────────────────────────────────────────────────────────────────────────
%% IEEE Access requires author biography. Replace placeholder text below.
%% ──────────────────────────────────────────────────────────────────────────
\begin{IEEEbiography}{Palak Srivastava}
received the B.Tech.\ degree in Computer Science and Engineering from
Siksha `O' Anusandhan University, ITER, Bhubaneswar, India, in 2023.
She is currently an independent researcher based in Hyderabad, India,
with research interests in decentralized systems, applied cryptography,
blockchain-based protocols, and peer-to-peer storage architectures.
This work investigates trustless storage markets on EVM-compatible
blockchains, focusing on browser-native encryption and on-chain
economic governance without privileged intermediaries.
\end{IEEEbiography}

\end{document}
